\subsection{Lista 1}

Vou passar a limpo exercĩcios do Varadhan.

\begin{exercise}[5.1]
	Construct a nonnegative martingale $X_n$ with $\Ex X_n = 1$ such that
	$\xi = \sup_n X_n$ is not integrable.
\end{exercise}
\begin{proof}
	The proof is basically done in the book. Just consider the
	Lebesgue masure on $[0,1]$ and the $\cF_n$ will be the sigma-algebras
	generated by the diadic sets of length $1/2^n$. Then define
	\[
		X_n = 2^n\id[\omega \in [0,2^{-n}]].
	\]
	It is easy to check it is a martingale. Furthermore,
	\[
		\int_{[0,1]} \xi dx = \sum_{j = 0}^{\infty} \int_{[2^{-j - 1}, 2^{j}]} X_j dx = \sum_{j = 0}^{\infty} 1/2 = \infty.
	\]
\end{proof}

\begin{exercise}[5.2]
	If $X_n$ is a martingale such that the differences $Y_n = X_{n} - X_{n-1}$ are all
	square integrable, show that for $n \neq m$, $\Ex[Y_n Y_m] = 0$. Therefore
	\[
		\Ex[X_n^2] = \Ex[X_0^2] + \sum_{j=1}^n \Ex[Y_j^2].
	\]
	If in addition, $\sup_n \Ex[X_n^2] < \infty$, then show there is a random variable $X$
	such that
	\[
		\lim_{n \to \infty} \Ex[|X - X_n|^2] = 0
	\]
\end{exercise}
\begin{proof}
	Without loss of generality, suppose $n > m$. Then, by simply conditioning on $\cF_{n-1}$
	\[
		\Ex[Y_n Y_m] = \Ex[(X_n - X_{n-1})Y_m] = \Ex[Y_m (\Ex[(X_n|\cF_{n-1}] - X_{n-1})] = 0.
	\]
	The other equality becomes obvious from the fact that $X_n = X_0 + \sum_{j=1}^n Y_j$.
	For the conclusion, note that together with the previous observation, it implies
	\[
		\sum_n \Ex Y_n^2 < \infty
	\]
	so that the $X_n$ form a Cauchy sequence in $L^2$. As $L^2$ is complete, there exists
	such $X$.
\end{proof}

\begin{exercise}[5.3]
	Show there exists a uniformly bounded in $L^1$ martingale $X_n$ but $X_n$ does
	not converge in $L^1$.
\end{exercise}
\begin{proof}
	The first exercise provides an example. It is clear that $X_n \to 0$ almost surely,
	but $|X_n|_{L^1} = 1$. This is an absurd, since any limit of $X_n$ in $L^1$ would have
	to be a.s equal to $\lim_{n\to\infty} X_n$.
\end{proof}

\begin{exercise}[5.4]
	Estabilish that a nonnegative martingale has an almost sure limit (which may not be in $L^1$).
	Show (I wont really do it) that we can assume, wlog that we are in the following situation:
	\[
		\Omega = \bigotimes_{j=1}^{\infty} \R, \quad \cF = \sigma(x_1, \dots, x_n),\quad X_j(\omega) = x_j.
	\]
	The existence of a $Q$ such that
	\[
		\frac{dQ}{dP}|_{\cF_n} = x_n
	\]
	is essentially a Kolgomorov's consistency theorem.
\end{exercise}
\begin{proof}
	Actually don't know how to do this one.
\end{proof}

\begin{exercise}[5.5]
	Show that $\tau(\omega) \equiv k$ is a stopping time for any constant $k$.
\end{exercise}
\begin{proof}
	Note that if $n < k$, then $\{\tau \leq n\} = \varnothing \in \cF_n$. Similarly, if $n \geq k$,
	$\{\tau \leq n\} = \Omega \in \cF_n$.
\end{proof}

\begin{exercise}[5.6]
	Show that if $\tau$ is a stopping time and $f : \N \to \N$ is a non-decreasing function
	with $f(t) \geq t$ for all $t$. Then $f(\tau)$ is also a stopping time.
\end{exercise}
\begin{proof}
	Just note that given $n$, there exists a maximal $m$ s.t $f(m) \leq n$ and by hypothesis, $m \leq n$. Then
	$\{f(\tau) \leq n\} = \{\tau \leq m\} \in \cF_m \subset \cF_n$.
\end{proof}

\begin{exercise}[5.7]
	Show that if $\tau_1$ and $\tau_2$ are stopping times, so is $\max(\tau_1, \tau_2)$ and $\min(\tau_1, \tau_2)$.
	In particular, $\tau$ is an increasing limit of bounded stopping times $\tau \land n$.
\end{exercise}
\begin{proof}
	Fix $n$, note that $\{\tau_1 \land \tau_2 \leq n\} = \{\tau_1 \leq n\} \cup \{\tau_1 \leq n\} \in \cF_n$.
	The same is true for the max, we just take intersection instead of union.
\end{proof}

\begin{exercise}[5.8]
	Verify that for a stopping time $\tau$, $\cF_\tau$ is a sub $\sigma$-field.
	If $\tau \equiv k$, then $\cF_\tau = \cF_k$. If $\tau_1 \leq \tau_2$ are stopping times,
	then $\cF_{\tau_1} \subset \cF_{\tau_2}$. Finally, if $\tau$ is a stopping time,
	then it is $\cF_\tau$ measurable.
\end{exercise}
\begin{proof}
	Let's show that $\cF_{\tau} \subset \cF$ is a $\sigma$-field. Clearly $\varnothing \in \cF_\{\tau\}$.
	Let $A \in \cF_{\tau}$, then for every $n$, $A \cap \{\tau \leq n\} \in \cF_{n}$,
	But then, $A^c \cup \{\tau > n\} \in \cF_n$, and therefore
	$A^c \cap \{\tau \leq n\} \in \cF_n$. Similarly, if $(A_n)_{n = 1}^{\infty}$ is
	a sequence of sets in $\cF_{\tau}$, then for any $m$
	\[
		\{\tau \leq m\} \cap \bigcup_{n\geq 1} A_n = \bigcup_{n\geq 1} \{\tau \leq m \} \cap A_n \in \cF_m.
	\]

	If $\tau \equiv k$. It is clear that $\cF_k \subset \cF_{\tau}$ since for any
	$A \in \cF_{k}$, $A \cap \{\tau \leq k\} = A \in \cF_k$.
	Similarly, if $A \in \cF_{\tau}$, then $A = A \cap {\tau = k} \in \cF_k$.

	If $\tau_1 \leq \tau_2$ and $A \in \cF_{\tau_1}$ then for any $m$
	$A \cap \{\tau_2 \leq m\} \subset (A \cap \{\tau_1 \leq m\}) \cap \{\tau_2 \leq m\}$.
	Note that $A \cap \{\tau_1 \leq m\} \in \cF_m$ and as $\{\tau \leq m\} \in \cF_m$,
	we have that their intersection is also in $\cF_m$.

	Finally, measurability is obvious since $\{\tau \leq n\} \in \cF_{\tau}$ for
	every $n$.
\end{proof}

\begin{exercise}[5.10]
	If $X_n$ is a submartingale and $\tau_1 \leq \tau_2 \leq M$ are bounded stopping times,
	we have that $\Ex[X_{\tau_2} | \tau_1] \geq X_{\tau_1}$. With the respective
	result for supermartingales.
\end{exercise}
\begin{proof}
	Decompose $X_n$ as $Y_n + A_n$ where $Y_n$ is a martingale and $A_n$ is strictly increasing and
	predictable. Now it is clear that, by linearity,
	\[
		\Ex[Y_{\tau_2} + A_{\tau_2} | \cF_{\tau_1}] = Y_{\tau_1} + \Ex[A_{\tau_2} | \cF_{\tau_1}] \geq Y_{\tau_1} + A_{\tau_1}.
	\]
	since $A_{\tau_2} \geq A_{\tau_1}$.

	The interesting case is when $X_n$ is a martingale. Then you may only assume $\tau \leq M$ since
	\[
		\Ex[X_m | \cF_{\tau_1}] = \Ex[\Ex[X_m | \cF_{\tau_2}] | \cF_{\tau_1}].
	\]
	Then we are left with proving that $\Ex[X_M | \cF_{\tau}] = X_{\tau}$, to do that just separate
	on the values of $\tau$.
\end{proof}

\begin{exercise}[5.11]
	Let $X_0 = 0$ and $X_n = \sum_{j \leq n} \xi_j$ where $\xi_j$ are u.i.i.d on $\{-1,1\}$.
	Let $\tau = \inf\{n : X_n = 1\}$. Then $\tau$ is a stopping time, $\Pr(\tau < \infty) = 1$,
	but $\tau$ is not bounded. $X_\tau = 1$ with probability $1$ but $\Ex[X_0] = 0$.
\end{exercise}
\begin{proof}
	I think the only thing to show on this exercise is that $\tau$ is a stopping time and
	$\Pr(\tau < \infty) = 1$. Note that $\{\tau = n\} = \{X_n = 1\} \cap \bigcap_{j \leq n} \{X_j \neq 1\} \in \cF_n$,
	so it is a stopping time. To show reccurence we can proceed with the usual
	Catalan counting or use markov chains arguments.
\end{proof}

\begin{exercise}[5.12]
	Let $\tau$ be an unbounded stopping time. For every $n$, $\tau \land n$ is a bounded stopping time
	and $\Ex X_{\tau \land n} = 0$. As $n \to \infty$, $\tau \land n \to \tau$. If we
	can estabilish uniform integrability, then we can pass to the limit. In particular,
	is $S(\omega) = \sup_{0 \leq n \leq \tau} |X_n(\omega)|$ is integrable, then $\sup_n |X_{\tau \land n}| < S$
	and so $\Ex X_\tau = 0$.
\end{exercise}

\begin{proof}
	Yes, if $Z_n$ converges almost surely to $Z$ and they are uniformly integrable, then $\Ex Z_n \to \Ex Z$.
	The way to prove this is with Egorov's theorem. If $\mu$ is a finite measure and $Z_n \to Z$ a.s,
	then for every $\eps$ there exists a set $A$ with $\mu(A^c) < \eps$ such that $Z_n \to Z$ uniformly in $A$.
	By using uniform integrability we can win over in $A^c$ and get usual convergence.
	Note that if $S$ is integrable, then for all $X_{\tau \land n} \ll S$, so the family is uniformly integrable.
	And we can apply the previous result.
\end{proof}


\begin{exercise} [5.14]
	Seja $\tau$ um stopping time e $\tau_k \to \tau$ uma sequência crescente
	de stopping times. Mostre que
	\[
		\cF_\tau = \sigma\bigg(\bigcup_k \cF_{\tau_k}\bigg)
	\]
\end{exercise}

Precisamos de algumas observações simples.

\begin{lemma}
	Sejam $\tau_1 \leq \tau_2$ stopping times. Então $\cF_{\tau_1}
		\subset \cF_{\tau_2}$.
\end{lemma}
\begin{proof}
	Tome $A \in \cF_{\tau_1}$. Por definição, para todo $n$ temos
	$A \cap \{\tau_1 \leq n\} \in \cF_n$. Como $\tau_1 \leq \tau_2$ e
	$\{\tau_2 \leq n\} \in \cF_n$,
	\[
		A \cap \{\tau_2 \leq n\} = (A \cap \{\tau_1 \leq n\}) \cap \{\tau_2
		\leq n\} \in \cF_n .
	\]
	Como isso vale para todo $n$, $A \in \cF_{\tau_2}$.
\end{proof}

\begin{obs}
	Seja $B \in \cF_n$ e $\tau$ um stopping time qualquer. Então vale que
	\[
		B \cap \{\tau = n\} \in \cF_{\tau} \quad \& \quad B \cap \{\tau >
		n\} \in \cF_{\tau}.
	\]
\end{obs}

\begin{proof}[Prova do Exercício]
	Segue imediatamente do primeiro lema que $\sigma(\bigcup_{k}
		\cF_{\tau_k}) \subset \cF_{\tau}$.
	Basta verificar a outra inclusão.
	Seja $A \in \cF_{\tau}$, separamos $A$ de acordo com o que acontece em $\tau$
	\[A = A \cap \{\tau = \infty\} \cup \bigg( \bigcup_{n} A \cap \{\tau
		\leq n\} \bigg).\]
	Vamos provar que cada pedacinho pertence a $\sigma(\bigcup_{k} \cF_{\tau_k})$.

	Por definição, $A \cap \{\tau \leq n\} \in \cF_n$. Como $\tau_k \to
		\tau$, segue que
	\[
		A \cap \{\tau \leq n\} = \bigcup_{k} [A \cap \{\tau \leq n\} \cap
			\{\tau_k = n\}].
	\]
	Pela observação, anterior, $A \cap \{\tau \leq n\} \cap \{\tau_k =
		n\} \in \cF_{\tau_k}$.
	Sendo união de conjuntos dessa forma, $A \cap \{\tau \leq n\} \in
		\sigma(\bigcup_k \cF_{\tau_k})$. Para $\{\tau = \infty\}$ temos basicamente a mesma coisa. Note que
	\[
		A \cap \{\tau = \infty\} = \bigcap_n A \cap \{\tau > n\}.
	\]
	Usando que $\tau_k \to \tau$,
	$A \cap \{\tau > n\} = \bigcup_k A \cap \{\tau > n\} \cap \{t_k > n\}$.
	Da mesma forma que antes, cada um dos termos pertence a
	$\cF_{\tau_k}$ e o resultado segue.
\end{proof}

\begin{exercise}
	Let $X_n = \sum_{j=1}^n \xi_j$ where each $\xi_j$ is a $\pm 1$ with probability $1/2$.
	Show that if $\tau$ is a stopping time with $\Ex \tau < \infty$, then
	$\sup_{1 \leq n \leq \tau} |X_n(\omega)|$ is square integrable and therefore $\Ex[X_\tau] = 0$.
	Hint: use the fact that $X_n^2 - n$ is a martingale.
\end{exercise}

\begin{proof}
	So, first we need to check $X_{n}^2 - n$ is a martingale, but this follows easily from
	a previous calculation and the fact that $\Ex \xi_i \xi_j = 0$ if $i \neq j$.
	Now let's use the previous exercise, note that considering $\tau \land n$ as a bounded stopping time,
	we get that $\Ex X_{\tau \land n}^2 - \tau \land n = 0$. So $\Ex X_{\tau \land n}^2 \leq \Ex \tau \land n \leq \Ex \tau$.
	It follows that $\Ex X_{\tau \land n}$ is uniformly bounded in $L^2$ and therefore uniformly
	integrable. So $0 = \lim_{n \to \infty }\Ex X_{\tau \land n} = \Ex X_{\tau}$.
\end{proof}

\begin{exercise}
	Show Doob's upcrossing inequality works for supermartingales, and chaging signs
	we get a downcrossing inequality for submartingales.
\end{exercise}

\begin{proof}
	The only thing we need to check is that for supermartingales, $\Ex D(\omega) \leq 0$. All the other
	inequalities are the same. By switching signs we get the result for submartingales.
\end{proof}

\begin{exercise}
	Let $X_n$ be a symmetric random walk on $\Z$ starting from $0$. Varadhan showed that for a fixed $R$ and $\lambda$ such that
	$\cos \lambda x \geq \cos \lambda R > 0$ for $x \in [-R, R]$ (we may take $\lambda \leq \pi/(2R)$) we have
	the inequality
	\[
		\Ex_x [e^{\sigma\tau_R}] \leq \frac{\cos \lambda x}{\cos \lambda R}
	\]
	where $\tau_R$ is the hitting time of $(-\infty, -R] \cup [R, \infty)$ and $\sigma = -\log \cos \lambda$.
	Can you show inequality above? What range is it valid in? Is $\Ex e^{\sigma \tau} < \infty$ for all $\sigma$?
\end{exercise}
\begin{proof}
	The way Varadhan proves this is by considering the martingale $e^{\sigma n} \cos(\lambda X_n)$.
	Now it is clear that for every $n$,
	\[
		\cos(\lambda x) = \Ex_x e^{\sigma \tau_R \land n} \cos(\lambda X_{n \land \tau}) =
		\Ex_x[\id[\tau > n] e^{\sigma n} \cos(\lambda X_n)] + \Ex_x[\id[\tau \leq n] e^{\sigma \tau} \cos(\lambda R)].
	\]
	If $\sigma < -\log \cos \lambda$, then the first term of the right hand side is bounded by
	$\Ex e^{\sigma n} \cos(\lambda X_n) \to 0$ when $n \to \infty$. As we know the second
	term converges to $\Ex e^{\sigma \tau} \cos(\lambda R)$ we get equality in this case.
	I should work a bit harder on the $\sigma = -\log \cos \lambda$ case.

	For the last observation, note that trivially $\Pr(X_1 = 1, X_2 = 0, \dots X_n = n\%2) = 2^{-n}$.
	So $\Ex e^{\sigma \tau} \geq 2^{-n}e^{\sigma n}$, in particular if $\sigma > \log 2$, as this is valid for all $n$,
	we already get $\Ex e^{\sigma \tau} = \infty$.
\end{proof}

